\hypertarget{criterion-7-institutional-support-and-financial-resources}{%
\chapter{Criterion 7: Institutional Support and Financial Resources}\label{criterion-7-institutional-support-and-financial-resources}}

\hypertarget{financial-governance-sustainability-c7}{%
\section{Financial Governance and Sustainability}\label{financial-governance-sustainability-c7}}

NUST manages financial resources through a centralized budgeting and monitoring framework involving academic departments, schools, and central administration. Annual allocations are based on academic requirements, enrolment, infrastructure needs, and research priorities. Recurring expenditure supports salaries, laboratory operations, maintenance, software and other routine requirements, while developmental expenditure supports laboratory equipment, computing infrastructure and facility enhancement. Budget utilization, procurement and audit are governed through established University procedures and applicable public procurement regulations.

SEECS demonstrates sustained financial capacity throughout the accreditation cycle. Table~\ref{tab:c7_financial_trends} summarizes the principal financial indicators; detailed account-head-wise budgeted and actual figures are provided in Annexure~\ref{annex:L}.

\begin{table}[H]
\centering
\caption{Key Financial Indicators, SEECS (Rupees in Million, Actual)}
\label{tab:c7_financial_trends}
\begingroup\footnotesize
\setlength{\tabcolsep}{4.5pt}
\renewcommand{\arraystretch}{1.12}
\begin{tabular}{@{}lrrrr@{}}
\toprule
\textbf{Indicator} & \textbf{2022--23} & \textbf{2023--24} & \textbf{2024--25} & \textbf{2025--26} \\
\midrule
Recurring expenditure & 486.408 & 508.979 & 622.386 & 731.382 \\
Development expenditure & 14.700 & 15.100 & 81.950 & 74.820 \\
Research project funding & 46.583 & 70.217 & 73.262 & 69.940 \\
Student tuition fee (self-funded) & 594.764 & 753.222 & 889.700 & 1107.514 \\
Scholarship support & 27.797 & 17.788 & 14.344 & 12.527 \\
\bottomrule
\end{tabular}
\endgroup
\end{table}

Recurring expenditure increased from Rs.~486.408~M in 2022--23 to Rs.~731.382~M in 2025--26, while developmental expenditure rose substantially in 2024--25 and 2025--26 due primarily to laboratory and infrastructure investment. In parallel, internally generated tuition income and externally funded research provide complementary revenue streams. These trends demonstrate continued institutional commitment and sufficient financial capacity to sustain the Software Engineering program.

\hypertarget{human-resource-academic-support-c7}{%
\section{Human Resource and Academic Support}\label{human-resource-academic-support-c7}}

A major component of recurring expenditure is the provision of qualified faculty and support staff. Actual expenditure on faculty and support-staff salaries increased from Rs.~461.508~M in 2022--23 to Rs.~632.311~M in 2025--26. NUST additionally supports employee benefits, faculty development, higher education, research incentives, conference participation and professional training. Faculty recruitment, qualifications, workload and current strength are reported in Criterion 5; duplicating these figures here is avoided to maintain a single authoritative record.

Academic delivery is further supported through centralized student services, information technology, learning-management systems, library and digital resources, and administrative support. These resources are funded through the University's recurring and developmental mechanisms and are available to the Software Engineering program as part of the broader SEECS/NUST institutional infrastructure.

\hypertarget{infrastructure-laboratory-support-c7}{%
\section{Infrastructure and Laboratory Support}\label{infrastructure-laboratory-support-c7}}

Institutional investment during the accreditation cycle has focused strongly on laboratory modernization, computing facilities and upkeep of existing infrastructure. The four-year actual expenditure shown in Table~\ref{tab:c7_lab_support} is derived from the detailed financial record in Annexure~\ref{annex:L}.

\begin{table}[H]
\centering
\caption{Four-Year Laboratory and Facility Support, SEECS (Rupees in Million, Actual)}
\label{tab:c7_lab_support}
\begingroup\footnotesize
\setlength{\tabcolsep}{8pt}
\renewcommand{\arraystretch}{1.12}
\begin{tabular}{@{}lr@{}}
\toprule
\textbf{Support Area} & \textbf{2022--23 to 2025--26} \\
\midrule
Laboratory equipment & 164.770 \\
Software purchase & 21.800 \\
Laboratory consumables & 3.000 \\
Maintenance of existing facilities & 153.888 \\
Repair and maintenance & 4.640 \\
\midrule
\textbf{Total} & \textbf{348.098} \\
\bottomrule
\end{tabular}
\endgroup
\end{table}

The increase in laboratory-equipment expenditure to Rs.~71.55~M in 2024--25 and Rs.~74.82~M in 2025--26 provides direct evidence of continued investment in modernization. Computing laboratories, networking, licensed software and associated facilities used by the program are described in detail under Criterion 6; the financial evidence here demonstrates the institutional capacity to establish, operate, maintain and upgrade those facilities.

\hypertarget{research-innovation-student-support-c7}{%
\section{Research, Innovation and Student Support}\label{research-innovation-student-support-c7}}

NUST supports research and innovation through institutional facilitation and externally sponsored projects. SEECS received approximately Rs.~260.0~M in research project funding over the four-year accreditation cycle, with annual funding remaining between Rs.~69.9~M and Rs.~73.3~M during the most recent three years. Faculty and students are supported through research infrastructure, project activities, consultancy and industry collaboration, prototype development, scholarly dissemination and participation in innovation activities.

Student financial support is provided through Board, GPA-based, HEC, MS-leading-to-PhD and need-based scholarship schemes. Cumulative scholarship support during 2022--23 to 2025--26 was approximately Rs.~72.46~M. Scheme-wise annual figures are retained in Annexure~\ref{annex:L} as primary evidence.

\hypertarget{criterion7-compliance-summary}{%
\section{Criterion 7 Compliance Summary}\label{criterion7-compliance-summary}}

\begin{table}[H]
\centering
\caption{Evidence Summary for PEC Criterion 7}
\label{tab:c7_evidence_summary}
\begingroup\footnotesize
\setlength{\tabcolsep}{7pt}
\renewcommand{\arraystretch}{1.08}
\begin{tabular}{@{}p{0.38\textwidth}p{0.50\textwidth}@{}}
\toprule
\textbf{Requirement} & \textbf{Primary Evidence} \\
\midrule
Financial sustainability & Table~\ref{tab:c7_financial_trends}; Annex~\ref{annex:L} \\
Continued financial commitment & Tables~\ref{tab:c7_financial_trends}--\ref{tab:c7_lab_support}; Annex L; Criterion 6 \\
Human-resource support & Salary expenditure; Criterion 5 faculty record \\
Research and innovation support & Research funding in Table~\ref{tab:c7_financial_trends}; project records \\
Student financial support & Scholarship expenditure; Annex~\ref{annex:L} \\
\bottomrule
\end{tabular}
\endgroup
\end{table}

The financial and institutional evidence demonstrates sustained support for academic operations, human resources, laboratories, research and students throughout the accreditation cycle. The upward trend in recurring expenditure, major developmental investment in the two most recent years, continued research funding and documented student support collectively establish that the Software Engineering program has adequate and sustainable institutional resources to meet the requirements of PEC Criterion~7.
